\par\medskip
\hypertarget{native-070-article053}{}\pdfbookmark[2]{53.}{article-053}
\begin{center}53.\end{center}
Situs puncti cuiuscunque in spatio commodissime per distantias a tribus planis sub angulis rectis se secantibus determinatur. Assumendo pro planorum uno planum eclipticae, designandoque per $z$ distantiam corporis coelestis ab hoc plano a parte boreali positive, ab australi negative sumendam, manifesto habebimus
\[
z=r'\operatorname{tang}\beta=r\sin\beta=r\sin i\sin u.
\]
Plana duo reliqua, quae per Solem quoque ducta supponemus, in sphaera coelesti circulos maximos proiicient, qui eclipticam sub angulis rectis secabunt, quorumque adeo poli in ipsa ecliptica iacebunt et $90^0$ ab invicem distabunt. Utriusque plani polum istum, a cuius parte distantiae positivae censentur, \textit{polum positivum} appellamus. Sint itaque $N$ et $N+90^0$ longitudines polorum positivorum, designenturque distantiae a planis quibus respondent respective per $x$, $y$. \SourceLeaf{071}Tunc facile perspicietur haberi
\[
\begin{aligned}
x&=r'\cos(\lambda-N)\\
&=r\cos\beta\cos(\lambda-\AscendingNode)\cos(N-\AscendingNode)\\
&\quad+r\cos\beta\sin(\lambda-\AscendingNode)\sin(N-\AscendingNode)\\[.4em]
y&=r'\sin(\lambda-N)\\
&=r\cos\beta\sin(\lambda-\AscendingNode)\cos(N-\AscendingNode)\\
&\quad-r\cos\beta\cos(\lambda-\AscendingNode)\sin(N-\AscendingNode),
\end{aligned}
\]
qui valores transeunt in
\[
\begin{aligned}
x&=r\cos(N-\AscendingNode)\cos u+r\cos i\sin(N-\AscendingNode)\sin u\\
y&=r\cos i\cos(N-\AscendingNode)\sin u-r\sin(N-\AscendingNode)\cos u.
\end{aligned}
\]
Quodsi itaque polus positivus plani ipsarum $x$ in ipso nodo ascendente collocatur, ut sit $N=\AscendingNode$, habebimus coordinatarum $x$, $y$, $z$ expressiones simplicissimas
\[
\begin{aligned}
x&=r\cos u\\
y&=r\cos i\sin u\\
z&=r\sin i\sin u.
\end{aligned}
\]
Si vero haec suppositio locum non habet, tamen formulae supra datae formam aeque fere commodam nanciscuntur per introductionem quatuor quantitatum auxiliarium $a$, $b$, $A$, $B$ ita determinatarum ut habeatur\HistoricalNoteEight{Handschriftliche Bemerkung, wozu vgl. Seite 69 Zeile 7:]
\[
\begin{aligned}
-\cos\AscendingNode'&=a\sin A\\
\cos i'\sin\AscendingNode'&=a\cos A\\
-\sin\AscendingNode'&=b\sin B\\
-\cos i'\cos\AscendingNode'&=b\cos B.
\end{aligned}
\]
}
\[
\begin{aligned}
\cos(N-\AscendingNode)&=a\sin A\\
\cos i\sin(N-\AscendingNode)&=a\cos A\\
-\sin(N-\AscendingNode)&=b\sin B\\
\cos i\cos(N-\AscendingNode)&=b\cos B
\end{aligned}
\]
(vid. art. 14, II). Manifesto tunc erit
\[
\begin{aligned}
x&=ra\sin(u+A)\\
y&=rb\sin(u+B)\\
z&=r\sin i\sin u.
\end{aligned}
\]
